\documentclass[10pt,a4paper]{article}
\usepackage[T1]{fontenc}\usepackage[utf8]{inputenc}\usepackage[english,italian]{babel}
\usepackage{lmodern,microtype}\usepackage[a4paper,margin=1.85cm,headheight=15pt]{geometry}
\usepackage{enumitem,tabularx,array,longtable,booktabs,xcolor,hyperref,xurl,fancyhdr,titlesec,framed}
\definecolor{ddtadark}{RGB}{28,38,52}\definecolor{ddtablue}{RGB}{42,88,135}\definecolor{ddtagray}{RGB}{244,246,248}\definecolor{ddtagreen}{RGB}{48,111,78}\definecolor{ddtaorange}{RGB}{148,91,25}
\hypersetup{colorlinks=true,linkcolor=ddtablue,urlcolor=ddtablue}\pagestyle{fancy}\fancyhf{}\lhead{DDTA - BA0-R}\rhead{REVIEW-ONLY}\cfoot{\thepage}
\titleformat{\section}{\large\bfseries\color{ddtadark}}{}{0pt}{}\setlist[itemize]{leftmargin=1.45em,itemsep=.18em,topsep=.2em}
\setlength{\parindent}{0pt}\setlength{\parskip}{.32em}\setlength{\emergencystretch}{2em}
\newenvironment{factbox}[1]{\def\FrameCommand{\fcolorbox{ddtablue}{ddtagray}}\MakeFramed{\advance\hsize-2\width\FrameRestore}\noindent\textbf{#1}\par\smallskip}{\endMakeFramed}
\newenvironment{challenge}[1]{\def\FrameCommand{\fcolorbox{ddtaorange}{ddtagray}}\MakeFramed{\advance\hsize-2\width\FrameRestore}\noindent\textbf{#1}\par\smallskip}{\endMakeFramed}

\begin{document}
\begin{center}
{\LARGE\bfseries DDTA - Scheda di lettura compilata}\par
{\large SRC-0046 - BA0-R systems-modeling prior art}\par
{\small REVIEW-ONLY: da approvare prima del caricamento nel repository.}\par
\end{center}
\begin{factbox}{Identita bibliografica e accesso}
\begin{tabularx}{\linewidth}{>{\bfseries}p{3.4cm}X}
Source ID & SRC-0046\\
Titolo & Survey of Candidate Model-Based Systems Engineering (MBSE) Methodologies, Revision B\\
Autori/ente & Jeff A. Estefan\\
Anno & 2008\\
Identificatore & INCOSE-TD-2007-003-02, Rev. B, May 23 2008\\
Accesso & public\_reference\_copy\\
Verification state & verified\_full\_text\_public\\
\end{tabularx}
\end{factbox}
\section{1. Research problem and contribution}
Provides a cursory survey of leading MBSE methodologies. It explicitly defines methodology as a collection of related processes, methods and tools and distinguishes those from the project environment. The report states that it is a survey, not a comparative assessment.
\section{2. Starting artifacts and generated representations}
Methodologies differ substantially. Harmony-SE derives functionality and states/modes and allocates them to physical architecture; OOSEM moves from stakeholder needs/requirements through logical and candidate allocated architectures; OPM integrates object/process/state concepts in one model. The diversity is evidence against assuming one established minimal ontology.
\section{3. Traceability, change, automation and human role}
The source discusses repositories, work products and methodology/tool integration but predates modern digital engineering. It is most useful for identifying second-ring sources rather than for current tool maturity or empirical effectiveness.
\section{4. Findings, limitations and relationship with DDTA}
The survey shows that structure/behavior/state/allocation recur in different methodologies, but with different conceptual organizations. That supports extracting semantic responsibilities before choosing metaclasses. It also yields bibliography candidates: Wymore, Harmony-SE, OOSEM, Vitech graphical representation relations, JPL State Analysis, Dori OPM and its reflective metamodel.
\section{Evidenza utile per BA0/BA1}
\begin{longtable}{p{1.2cm}p{2.0cm}p{4.0cm}p{8.0cm}}
\toprule
Pag. & Ruolo & Sezione & Interpretazione DDTA\\
\midrule
\endhead
0 & method & 1.1 Purpose / 1.2 Scope & Use the source as a map/bibliography seed, not as evidence that one method is superior.\\
\addlinespace[.3em]
2 & foundation & 2.1 PMTE definitions & Environment/context is a separate semantic concern from the modeled system elements.\\
\addlinespace[.3em]
15 & comparison & 3.1 Harmony-SE & Behavior, state/mode, structure and allocation recur as distinct responsibilities.\\
\addlinespace[.3em]
18 & comparison & 3.2 OOSEM & A neutral representation should support transition between logical and allocated views without equating them.\\
\addlinespace[.3em]
42 & challenge & 3.6 OPM & Alternative modeling kernels can be much smaller than SysML while preserving multiple semantic dimensions.\\
\addlinespace[.3em]
59 & bibliography & References & Bibliography mining is the primary value of this source for BA0-R.\\
\addlinespace[.3em]
64 & comparison & Appendix OPM concepts & State and transformation semantics recur even in a very different modeling paradigm.\\
\addlinespace[.3em]
\bottomrule
\end{longtable}
\section{Bibliography mining / follow-up}
\begin{itemize}
\item Dori 2002 Object-Process Methodology
\item Reinhartz-Berger \& Dori 2005 reflective OPM metamodel
\item Long 2002 Relationships between Common Graphical Representations in Systems Engineering
\item OOSEM/Harmony-SE foundational references
\item JPL State Analysis sources
\end{itemize}
\begin{challenge}{Governance note}
Questa scheda e un ausilio di review. Le future citazioni di tesi devono risolversi alla fonte registrata e a una posizione esatta nell'excerpt ledger approvato. Nessun concetto esterno diventa automaticamente un BAE DDTA.
\end{challenge}
\end{document}
